% =========================================================
\section{Semaine 5 : Interpolation}
% =========================================================

\begin{explication}[Idée principale]
Construire une fonction continue $p$ telle que $p(x_i) = y_i$ pour $i = 0, \dots, n$.
\end{explication}

\subsection{Interpolation polynomiale}

Étant donnés des points $(x_i, y_i)$ pour $i = 0, \dots, n$, on aimerait construire $p$, un polynôme de degré $n$, tel que :
$$p(x_i) = y_i, \quad \forall i = 0, \dots, n$$

Posons $p(x) = a_0 + a_1x + a_2x^2 + \dots + a_nx^n$.
La condition $p(x_i) = y_i$ équivaut à :
$$a_0 + a_1x_i + \dots + a_nx_i^n = y_i, \quad \forall i = 0, \dots, n$$

Ce qui nous donne le système matriciel suivant :
$$
\begin{bmatrix}
1 & x_0 & x_0^2 & \dots & x_0^n \\
1 & x_1 & x_1^2 & \dots & x_1^n \\
\vdots & \vdots & \vdots & \ddots & \vdots \\
1 & x_n & x_n^2 & \dots & x_n^n
\end{bmatrix}
\begin{bmatrix}
a_0 \\
a_1 \\
\vdots \\
a_n
\end{bmatrix}
=
\begin{bmatrix}
y_0 \\
y_1 \\
\vdots \\
y_n
\end{bmatrix}
$$
Où la matrice du système est $V$, la \textbf{matrice de Vandermonde}.

Si les $x_i$ sont tous distincts, il existe un unique polynôme de degré $n$ passant par les points $(x_i, y_i)$, $i=0,\dots,n$. Cela implique que le système linéaire est inversible.

\begin{remarque}[Problème avec la matrice de Vandermonde]
\textbf{Mais !} Il faut résoudre un système de taille $(n+1) \times (n+1)$, ce qui est \textbf{coûteux}. \\
De plus, le conditionnement de $V$, noté $\kappa(V)$, est très grand. Par conséquent, une petite erreur sur $y$ entraîne une grande différence sur le résultat.

En effet, si $Va = y$ et $V\tilde{a} = \tilde{y}$, on a :
$$\frac{\|a - \tilde{a}\|}{\|a\|} \le \kappa(V) \frac{\|y - \tilde{y}\|}{\|y\|}$$
Si $\kappa(V)$ est grand, même si l'erreur relative $\frac{\|y - \tilde{y}\|}{\|y\|} \ll 1$, l'erreur sur les coefficients $\frac{\|a - \tilde{a}\|}{\|a\|}$ sera potentiellement très grande.
\end{remarque}

\subsection{Base de Lagrange}

Au lieu de chercher $p$ dans la base canonique $\{1, x, \dots, x^n\}$, on va chercher $p$ dans une base $\{L_i\}_{i=0}^n$.

On aimerait que $L_i$ soit un polynôme de degré $n$ tel que :
$$
L_i(x_j) = 
\begin{cases} 
1 & \text{si } i = j \\ 
0 & \text{sinon} 
\end{cases}
$$
avec $x_i$, $i=0,\dots,n$, les nœuds d'interpolation.

\textbf{Pourquoi ?} \\
Si l'on écrit $p(x) = \sum_{i=0}^n \alpha_i L_i(x)$, alors :
$$p(x_j) = \sum_{i=0}^n \alpha_i L_i(x_j) = \alpha_j$$
Donc, pour assurer $p(x_j) = y_j$ pour $j=0,\dots,n$, il suffit de poser :
$$p(x) = \sum_{i=0}^n y_i L_i(x)$$

\subsubsection{Construction de la base de Lagrange}

\begin{definition}[Polynômes de Lagrange]
On définit les polynômes de Lagrange $L_i(x)$ par :
$$L_i(x) = \prod_{\substack{j=0 \\ j \neq i}}^n \frac{(x - x_j)}{(x_i - x_j)}$$
\end{definition}

Alors $L_i$ est bien un polynôme de degré $n$.
Il reste à montrer que c'est bien une base de $\mathbb{P}_n$ (l'ensemble des polynômes de degré $\le n$).

Pour ceci, puisque $\mathbb{P}_n$ est un espace vectoriel de dimension $(n+1)$, et qu'on a $(n+1)$ fonctions $L_i$, il suffit de montrer que c'est une famille libre.

\begin{explication}[Indépendance linéaire]
Supposons qu'il existe des coefficients $c_0, \dots, c_n$ tels que $\sum_{i=0}^n c_i L_i(x) = 0, \forall x$. \\
En particulier, si on évalue en $x = x_j$ :
$$0 = \sum_{i=0}^n c_i L_i(x_j) = c_j$$
Donc $c_j = 0, \forall j$. La famille est donc bien une famille libre.
\end{explication}

La famille $\{L_i\}_{i=0}^n$ est appelée la \textbf{base de Lagrange} associée aux points $x_i$ ($i=0,\dots,n$).

\textbf{En résumé, pour construire $p$ tel que $p(x_i) = y_i$ :}
\begin{enumerate}
    \item On construit la base de Lagrange $\{L_i\}_{i=0}^n$ associée aux points $\{x_i\}_{i=0}^n$.
    \item On définit le polynôme interpolant : $p(x) = \sum_{i=0}^n y_i L_i(x)$.
\end{enumerate}

\subsection{Interpolation de fonction}

Étant donnée une fonction $f : [a, b] \to \mathbb{R}$ continue, et des points $x_i \in [a, b]$, on aimerait construire $p_n$, un polynôme de degré $n$, tel que :
$$p_n(x_i) = f(x_i), \quad \forall i = 0, \dots, n$$

\textbf{Pourquoi fait-on cela ?}
\begin{itemize}
    \item Évaluation plus rapide de $f$
    \item Représentation graphique de $f$
    \item Intégration numérique
    \item Dérivation numérique
\end{itemize}

Comme vu précédemment, étant donnés les $x_i$, on pose :
$$p_n(x) = \sum_{i=0}^n f(x_i) L_i(x)$$
qui est l'interpolant de degré $n$ de $f$ aux points $x_i$ ($i=0,\dots,n$).

\textbf{Question :} Est-ce que $p_n$ est une bonne approximation de $f$ ? Est-ce que $p_n \to f$ quand $n \to \infty$ ?

\begin{theoreme}[Erreur d'interpolation (nœuds quelconque)]
Soit $f : [a,b] \to \mathbb{R}$ une fonction de classe $C^{n+1}$. Soient $x_0, \dots, x_n \in [a,b]$. \\
Soit $p_n$ le polynôme d'interpolation de $f$ aux points $x_i$. Alors :
$$\max_{x\in[a,b]} |f(x) - p_n(x)| \le \frac{1}{(n+1)!} \max_{\xi\in[a,b]} |f^{(n+1)}(\xi)| \max_{x\in[a,b]} \prod_{i=0}^n |x - x_i|$$
\end{theoreme}

\begin{theoreme}[Erreur d'interpolation (nœuds équidistribués)]
Soit $f : [a,b] \to \mathbb{R}$ une fonction de classe $C^{n+1}$. \\
Soient $x_0, \dots, x_n \in [a,b]$ équidistribués (i.e., avec un pas $h = \frac{b-a}{n}$). Alors :
$$\max_{x\in[a,b]} |f(x) - p_n(x)| \le \frac{1}{(n+1)!} \left(\frac{b-a}{n}\right)^{n+1} \max_{\xi\in[a,b]} |f^{(n+1)}(\xi)|$$
\end{theoreme}

\begin{remarque}[Instabilités (Phénomène de Runge)]
Pour l'interpolation polynomiale globale, quand $n \to \infty$, il y a des \textbf{instabilités}. Il faut que la fonction $f$ soit extrêmement régulière pour que la série d'interpolation converge.
\end{remarque}

\subsection{Interpolation par intervalles (par morceaux / Splines)}

Pour contourner le problème des grandes valeurs de $n$, on considère $(N+1)$ points :
$$a = x_0 < x_1 < \dots < x_N = b$$

On approxime $f$ par morceaux. Sur chaque sous-intervalle $[x_i, x_{i+1}]$ (pour $i=0,\dots,N-1$), on considère $n$ points équidistribués :
$$x_i = y_{i,0} < y_{i,1} < \dots < y_{i,n} = x_{i+1}$$

On interpole localement par un polynôme de degré $n$ passant par les points $(y_{i,j}, f(y_{i,j}))$. \\
En faisant cela, on approxime $f$ par une fonction $f_h$ qui est :
\begin{itemize}
    \item L'interpolant de degré $n$ par intervalles de $f$.
    \item Continue sur $[a,b]$.
    \item Un polynôme de degré $n$ strictement sur chaque sous-intervalle $[x_i, x_{i+1}]$.
\end{itemize}

\begin{theoreme}[Erreur d'interpolation par morceaux]
Soit $f : [a,b] \to \mathbb{R}$ une fonction $C^{n+1}$ et $f_h$ son interpolant de degré $n$ par intervalles. \\
Posons $h = \max_i |x_{i+1} - x_i|$. Alors il existe une constante $C$ (indépendante de $h$) telle que :
$$\max_{x\in[a,b]} |f(x) - f_h(x)| \le C h^{n+1}$$
En particulier, si les $x_i$ sont équidistribués avec un pas $h$, alors cette borne s'applique.
\end{theoreme}

\textbf{En pratique :} on prend $n$ petit ($n = 1, 2, 3, 4$) et $N$ grand (beaucoup de sous-intervalles).

\begin{explication}[Démonstration de la borne d'erreur par morceaux]
En utilisant le théorème précédent appliqué localement sur le sous-intervalle $[x_i, x_{i+1}]$, on a :
$$\max_{x\in[x_i, x_{i+1}]} |f(x) - f_h(x)| \le \frac{1}{(n+1)!} \left(\frac{x_{i+1}-x_i}{n}\right)^{n+1} \max_{\xi \in [x_i, x_{i+1}]} |f^{(n+1)}(\xi)|$$
On majore $x_{i+1}-x_i$ par le pas global $h$, et le maximum de la dérivée locale par le maximum sur tout le domaine $[a,b]$ :
$$\le \frac{1}{(n+1)! \, n^{n+1}} \max_{\xi\in[a,b]} |f^{(n+1)}(\xi)| \, h^{n+1}$$
Ce qui nous donne bien l'estimation globale :
$$\max_{x\in[a,b]} |f(x) - f_h(x)| \le C h^{n+1}$$
avec $C = \frac{1}{(n+1)! \, n^{n+1}} \max_{\xi\in[a,b]} |f^{(n+1)}(\xi)|$.
\end{explication}